\section{Offline-First Advantages}
\label{sec:offline}

The capsule architecture is offline-first by construction, not as an
afterthought. This produces several advantages beyond raw latency.

\subsection{Mobile and Edge Deployments}

Mobile applications frequently operate in degraded network conditions:
subway tunnels, airplane mode, rural areas with intermittent connectivity.
Firebase's offline mode handles this with a cache, but the cache is a
subset of the data, and complex queries may fail. The capsule model has
no degradation: the full database is on the device, and all operations
work identically regardless of network state.

\subsection{Data Residency}

For applications subject to data residency regulations (GDPR, HIPAA,
financial regulations), the capsule model simplifies compliance: the
data resides on the user's device. There is no question of which data
center holds the data, which jurisdiction's courts can compel access,
or whether cross-border data transfer provisions apply. The user's
device is the primary storage location.

\begin{proposition}[Data Residency by Construction]
In the capsule model, the user's data resides on the user's device by
default. Storage providers hold encrypted blobs for durability, but
the plaintext data never leaves the user's device (or TEE). This
satisfies data residency requirements without geographic provisioning
of cloud infrastructure, because the ``residence'' is the user's
device, which is by definition in the user's jurisdiction.
\end{proposition}

\subsection{Latency-Sensitive Applications}

Applications that require sub-millisecond read latency (real-time games,
audio/video processing, local AI inference with context retrieval) cannot
tolerate network round-trips. The capsule model enables these applications
by making the database local: a retrieval-augmented generation (RAG)
pipeline that queries a local vector index over a local SQLite database
operates with microsecond latency, compared to millisecond latency for
a cloud-hosted vector database.

